\documentclass[10pt,a4paper]{amsart}
\usepackage{amsmath,amssymb,amsthm,graphicx}
\usepackage{verbatim}
\newenvironment{reference}{\comment}{\endcomment}
%\newenvironment{reference}{}{}
\newenvironment{slogan}{\comment}{\endcomment}
\newenvironment{history}{\comment}{\endcomment}

% For commutative diagrams we use Xy-pic
\usepackage[all]{xy}

% We use 2cell for 2-commutative diagrams.
\xyoption{2cell}
% Theorem environments.
%
\theoremstyle{plain}
\newtheorem{theorem}[subsection]{Theorem}
\newtheorem{proposition}[subsection]{Proposition}
\newtheorem{lemma}[subsection]{Lemma}

\theoremstyle{definition}
\newtheorem{definition}[subsection]{Definition}
\newtheorem{example}[subsection]{Example}
\newtheorem{exercise}[subsection]{Exercise}
\newtheorem{situation}[subsection]{Situation}

\theoremstyle{remark}
\newtheorem{remark}[subsection]{Remark}
\newtheorem{remarks}[subsection]{Remarks}

\numberwithin{equation}{subsection}

% Macros
%
\def\lim{\mathop{\mathrm{lim}}\nolimits}
\def\colim{\mathop{\mathrm{colim}}\nolimits}
\def\Spec{\mathop{\mathrm{Spec}}}
\def\Hom{\mathop{\mathrm{Hom}}\nolimits}
\def\Ext{\mathop{\mathrm{Ext}}\nolimits}
\def\SheafHom{\mathop{\mathcal{H}\!\mathit{om}}\nolimits}
\def\SheafExt{\mathop{\mathcal{E}\!\mathit{xt}}\nolimits}
\def\Sch{\mathit{Sch}}
\def\Mor{\mathop{\mathrm{Mor}}\nolimits}
\def\Ob{\mathop{\mathrm{Ob}}\nolimits}
\def\Sh{\mathop{\mathit{Sh}}\nolimits}
\def\NL{\mathop{N\!L}\nolimits}
\def\CH{\mathop{\mathrm{CH}}\nolimits}
\def\proetale{{pro\text{-}\acute{e}tale}}
\def\etale{{\acute{e}tale}}
\def\QCoh{\mathit{QCoh}}
\def\Ker{\mathop{\mathrm{Ker}}}
\def\Im{\mathop{\mathrm{Im}}}
\def\Coker{\mathop{\mathrm{Coker}}}
\def\Coim{\mathop{\mathrm{Coim}}}

% Boxtimes
%
\DeclareMathSymbol{\boxtimes}{\mathbin}{AMSa}{"02}

%
% Macros for moduli stacks/spaces
%
\def\QCohstack{\mathcal{QC}\!\mathit{oh}}
\def\Cohstack{\mathcal{C}\!\mathit{oh}}
\def\Spacesstack{\mathcal{S}\!\mathit{paces}}
\def\Quotfunctor{\mathrm{Quot}}
\def\Hilbfunctor{\mathrm{Hilb}}
\def\Curvesstack{\mathcal{C}\!\mathit{urves}}
\def\Polarizedstack{\mathcal{P}\!\mathit{olarized}}
\def\Complexesstack{\mathcal{C}\!\mathit{omplexes}}
% \Pic is the operator that assigns to X its picard group, usage \Pic(X)
% \Picardstack_{X/B} denotes the Picard stack of X over B
% \Picardfunctor_{X/B} denotes the Picard functor of X over B
\def\Pic{\mathop{\mathrm{Pic}}\nolimits}
\def\Picardstack{\mathcal{P}\!\mathit{ic}}
\def\Picardfunctor{\mathrm{Pic}}
\def\Deformationcategory{\mathcal{D}\!\mathit{ef}}
\usepackage{cleveref}
\crefname{originallemma}{original lemma}{original lemmas}
\Crefname{originallemma}{Original lemma}{Original lemmas}
\makeatletter
\newlabel{dualizing-lemma-base-change-dualizing-over-field}{{0E0P}{source}{}{sourcepointer-1}{}}
\newlabel{dualizing-lemma-base-change-dualizing-over-field@cref}{{[originallemma][0][]0E0P}{[1][0][]source}{}{}{}}
\makeatother
\title{The Gorenstein locus of a flat locally finite type morphism commutes with arbitrary base change}
\author{The Stacks Project Authors}
\date{stacks-project:duality.tex:local-snapshot-20261012; GNU FDL 1.2 or later}
\begin{document}
\maketitle
\noindent\textit{Editorial note.} All delivered mathematical text is the author's own, copied from the pinned Stacks Project chapter file duality.tex; no statement or proof step was rewritten, reordered or removed.\par\medskip
\noindent\textit{Editorial note.} Preamble provenance: the pinned chapter file opens with an input of the shared preamble and carries no preamble text of its own, so the project's single Stacks-wide preamble staged at sources/stacks-duality/preamble.tex is excerpted verbatim under the recipe fixed by LANE\_BRIEF 22.1, as three spans, lines 8-13 (verbatim package and the reference/slogan/history environments), lines 14-18 (Xy-pic setup) and lines 154-217 (the native theorem declarations and mathematical macros). Nothing was added, removed or redefined; the wrapper class and the amsmath/amssymb/amsthm/graphicx base packages are supplied by the assembler.\par\medskip
\noindent\textit{Editorial note.} The counted statement and proof contain no cross reference, no citation and no included file, so the delivered body is self contained.\par\medskip
\noindent\textit{Editorial note.} Canonical location of the counted proposition: Stacks tag 0E0Q (sources/stacks-duality/pinned-original-source.tex line 6683).\par\medskip
\noindent\textit{Allowed standard prerequisites (editorial).}\begin{itemize}
\item Base change of a dualizing complex over a field. Dualizing Complexes, Lemma labelled dualizing-lemma-base-change-dualizing-over-field, cited at proof lines 6710-6712. Shows the base changed complex is again dualizing.
\item Flatness of the localised ring map and invertibility of the localised dualizing complex. The flat local homomorphism A\_p -> (A\_K)\_\{p\_K\} and the invertibility criterion used at proof lines 6714-6719. Relates Gorensteinness of the two local rings.
\item Local Noetherianity of the fibres of a flat locally finite type morphism. The observation at proof lines 6690-6694. Makes the statement meaningful and stable under base change.
\end{itemize}\medskip
\ifdefined\hypertarget\hypertarget{sourcepointer-1}{}\fi
\noindent\textit{Source reference.} Original lemma 0E0P: sources/stacks-chow/tags.txt, source lines 17544--17544. Cross-chapter standard prerequisite located in the project tag list staged in the local snapshot; invoked at proof line 6709 of the counted proof. It is retained only as a navigation pointer; every load-bearing step is inside the delivered span.\par\medskip
\begin{lemma}
\label{lemma-flat-lft-base-change-gorenstein}
Let $f : X \to Y$ be a morphism of schemes which is flat and
locally of finite type. Then formation of the set
$\{x \in X \mid f\text{ is Gorenstein at }x\}$
commutes with arbitrary base change.
\end{lemma}

\begin{proof}
The assumption implies any fibre of $f$ is locally of finite type
over a field and hence locally Noetherian and the same is true for
any base change. Thus the statement makes sense. Looking at
fibres we reduce to the following problem: let $X$ be a scheme
locally of finite type over a field $k$,
let $K/k$ be a field extension, and
let $x_K \in X_K$ be a point with image $x \in X$.
Problem: show that $\mathcal{O}_{X_K, x_K}$ is Gorenstein if and only if
$\mathcal{O}_{X, x}$ is Gorenstein.

\medskip\noindent
The problem can be solved using a bit of algebra as follows.
Choose an affine open $\Spec(A) \subset X$ containing $x$.
Say $x$ corresponds to $\mathfrak p \subset A$.
With $A_K = A \otimes_k K$ we see that $\Spec(A_K) \subset X_K$
contains $x_K$. Say $x_K$ corresponds to $\mathfrak p_K \subset A_K$.
Let $\omega_A^\bullet$ be a dualizing complex for $A$.
By Dualizing Complexes, Lemma
\ref{dualizing-lemma-base-change-dualizing-over-field}
$\omega_A^\bullet \otimes_A A_K$ is a dualizing complex for $A_K$.
Now we are done because
$A_\mathfrak p \to (A_K)_{\mathfrak p_K}$ is a flat local
homomorphism of Noetherian rings and hence
$(\omega_A^\bullet)_\mathfrak p$ is an invertible object
of $D(A_\mathfrak p)$ if and only if
$(\omega_A^\bullet)_\mathfrak p \otimes_{A_\mathfrak p} (A_K)_{\mathfrak p_K}$
is an invertible object of $D((A_K)_{\mathfrak p_K})$.
Some details omitted; hint: look at cohomology modules.
\end{proof}

\end{document}
